\documentclass[10pt,a4paper]{amsart}
\usepackage{amsmath,amssymb,amsthm,graphicx}
\usepackage{verbatim}
\newenvironment{reference}{\comment}{\endcomment}
%\newenvironment{reference}{}{}
\newenvironment{slogan}{\comment}{\endcomment}
\newenvironment{history}{\comment}{\endcomment}
% For commutative diagrams we use Xy-pic
\usepackage[all]{xy}

% We use 2cell for 2-commutative diagrams.
\xyoption{2cell}
% Theorem environments.
%
\theoremstyle{plain}
\newtheorem{theorem}[subsection]{Theorem}
\newtheorem{proposition}[subsection]{Proposition}
\newtheorem{lemma}[subsection]{Lemma}

\theoremstyle{definition}
\newtheorem{definition}[subsection]{Definition}
\newtheorem{example}[subsection]{Example}
\newtheorem{exercise}[subsection]{Exercise}
\newtheorem{situation}[subsection]{Situation}

\theoremstyle{remark}
\newtheorem{remark}[subsection]{Remark}
\newtheorem{remarks}[subsection]{Remarks}

\numberwithin{equation}{subsection}

% Macros
%
\def\lim{\mathop{\mathrm{lim}}\nolimits}
\def\colim{\mathop{\mathrm{colim}}\nolimits}
\def\Spec{\mathop{\mathrm{Spec}}}
\def\Hom{\mathop{\mathrm{Hom}}\nolimits}
\def\Ext{\mathop{\mathrm{Ext}}\nolimits}
\def\SheafHom{\mathop{\mathcal{H}\!\mathit{om}}\nolimits}
\def\SheafExt{\mathop{\mathcal{E}\!\mathit{xt}}\nolimits}
\def\Sch{\mathit{Sch}}
\def\Mor{\mathop{\mathrm{Mor}}\nolimits}
\def\Ob{\mathop{\mathrm{Ob}}\nolimits}
\def\Sh{\mathop{\mathit{Sh}}\nolimits}
\def\NL{\mathop{N\!L}\nolimits}
\def\CH{\mathop{\mathrm{CH}}\nolimits}
\def\proetale{{pro\text{-}\acute{e}tale}}
\def\etale{{\acute{e}tale}}
\def\QCoh{\mathit{QCoh}}
\def\Ker{\mathop{\mathrm{Ker}}}
\def\Im{\mathop{\mathrm{Im}}}
\def\Coker{\mathop{\mathrm{Coker}}}
\def\Coim{\mathop{\mathrm{Coim}}}

% Boxtimes
%
\DeclareMathSymbol{\boxtimes}{\mathbin}{AMSa}{"02}

%
% Macros for moduli stacks/spaces
%
\def\QCohstack{\mathcal{QC}\!\mathit{oh}}
\def\Cohstack{\mathcal{C}\!\mathit{oh}}
\def\Spacesstack{\mathcal{S}\!\mathit{paces}}
\def\Quotfunctor{\mathrm{Quot}}
\def\Hilbfunctor{\mathrm{Hilb}}
\def\Curvesstack{\mathcal{C}\!\mathit{urves}}
\def\Polarizedstack{\mathcal{P}\!\mathit{olarized}}
\def\Complexesstack{\mathcal{C}\!\mathit{omplexes}}
% \Pic is the operator that assigns to X its picard group, usage \Pic(X)
% \Picardstack_{X/B} denotes the Picard stack of X over B
% \Picardfunctor_{X/B} denotes the Picard functor of X over B
\def\Pic{\mathop{\mathrm{Pic}}\nolimits}
\def\Picardstack{\mathcal{P}\!\mathit{ic}}
\def\Picardfunctor{\mathrm{Pic}}
\def\Deformationcategory{\mathcal{D}\!\mathit{ef}}
\usepackage{cleveref}
\crefname{originallemma}{original lemma}{original lemmas}
\Crefname{originallemma}{Original lemma}{Original lemmas}
\makeatletter
\newlabel{lemma-check-quasi-compact-covering}{{lemma-check-quasi-compact-covering}{source}{}{sourcepointer-1}{}}
\newlabel{lemma-check-quasi-compact-covering@cref}{{[originallemma][0][]lemma-check-quasi-compact-covering}{[1][0][]source}{}{}{}}
\makeatother
\crefname{originalreference}{original reference}{original references}
\Crefname{originalreference}{Original reference}{Original references}
\makeatletter
\newlabel{stacks-properties-lemma-characterize-surjective}{{stacks-properties-lemma-characterize-surjective}{source}{}{sourcepointer-2}{}}
\newlabel{stacks-properties-lemma-characterize-surjective@cref}{{[originalreference][0][]stacks-properties-lemma-characterize-surjective}{[1][0][]source}{}{}{}}
\makeatother
\title{Decency descends along a quasi-compact surjective morphism}
\author{The Stacks Project Authors}
\date{stacks-project:stacks-morphisms.tex:local-snapshot-20261012; GNU FDL 1.2 or later}
\begin{document}
\maketitle
\noindent\textit{Editorial note.} All delivered mathematical text is the author's own, copied from the pinned Stacks Project chapter source sources/stacks-algebraic-stacks-morphisms/pinned-original-source.tex; no statement or proof step was rewritten, reordered or removed.\par\medskip
\noindent\textit{Editorial note.} Preamble provenance: the Stacks Project keeps one single shared preamble for every chapter; this card takes it directly from the staged sources/stacks-algebraic-stacks-morphisms/preamble.tex as three verbatim spans (lines 8-12, 14-18 and 154-217), so no derived file was written and nothing was added or redefined. The wrapper class and the base packages are supplied by the assembler.\par\medskip
\noindent\textit{Editorial note.} Bibliography provenance: neither the counted statement nor its proof contains a citation (this lane rescans the whole span before writing the card), so the card carries no bibliography and no entry was reconstructed.\par\medskip
\noindent\textit{Editorial note.} Selection provenance: this card was opened under the file-based control contract (CONTROL.md), with the mechanical pre-card gate run on its statement and proof interval before the card was written.\par\medskip
\noindent\textit{Allowed standard prerequisites (editorial).}\begin{itemize}
\item The characterisation of surjectivity of a morphism of algebraic stacks through the surjectivity on points, and the criterion for decency through the quasi-compactness of a fibre product along a presentation. Invoked inside the counted proof (source lines 10473-10480) as pointers to results of the chapter on properties of stacks and of this chapter. Reduces the assertion to the quasi-compactness of a fibre product over a field valued point.
\item The quasi-compactness of the morphism from a decent stack to a base along a field valued point, and the stability of quasi-compactness of a fibre product under such a point. Used inside the counted proof (source lines 10481-10489). Provides the two quasi-compact spaces whose fibre product computes the required one.
\end{itemize}\medskip
\ifdefined\hypertarget\hypertarget{sourcepointer-1}{}\fi
\noindent\textit{Source reference.} Original lemma lemma-check-quasi-compact-covering: sources/stacks-algebraic-stacks-morphisms/pinned-original-source.tex, source lines 2243--2243. Navigation pointer only: the counted text refers to the Stacks item lemma-check-quasi-compact-covering, which belongs to another chapter of the project and is not part of this card. The pointer carries the target's own label as its number because this chapter's staged snapshot does not contain the other chapter's tag table; its content is not reproduced here.\par\medskip
\ifdefined\hypertarget\hypertarget{sourcepointer-2}{}\fi
\noindent\textit{Source reference.} Original reference stacks-properties-lemma-characterize-surjective: sources/stacks-algebraic-stacks-morphisms/pinned-original-source.tex, source lines 4141--4141. Navigation pointer only: the counted text refers to the Stacks item stacks-properties-lemma-characterize-surjective, which belongs to another chapter of the project and is not part of this card. The pointer carries the target's own label as its number because this chapter's staged snapshot does not contain the other chapter's tag table; its content is not reproduced here.\par\medskip
\begin{lemma}
\label{lemma-qc-decent}
Let $f : \mathcal{X} \to \mathcal{Y}$ be a morphism of algebraic stacks.
If $f$ is quasi-compact and surjective and $\mathcal{X}$ is decent,
then $\mathcal{Y}$ is decent.
\end{lemma}

\begin{proof}
Let $x : \Spec(k) \to \mathcal{X}$ be a morphism where $k$ is a field
and denote $y = f \circ x$. Since $f$ is surjective, every point
of $|\mathcal{Y}|$ arises in this manner, see Properties of Stacks,
Lemma \ref{stacks-properties-lemma-characterize-surjective}.
Consider an affine scheme $T$ and morphism
$T \to \mathcal{Y}$. It suffices to show that
$T \times_{\mathcal{Y}, y} \Spec(k)$ is quasi-compact, see
Lemma \ref{lemma-check-quasi-compact-covering}.
We have
$$
(T \times_{\mathcal{Y}} \mathcal{X}) \times_{\mathcal{X}, x} \Spec(k) =
T \times_{\mathcal{Y}, y} \Spec(k)
$$
The morphism $T \times_{\mathcal{Y}} \mathcal{X} \to T$ is quasi-compact
hence $T \times_\mathcal{Y} \mathcal{X}$ is quasi-compact. Since $x$ is
a quasi-compact morphism as $\mathcal{X}$ is decent
we see that the displayed fibre product is quasi-compact.
\end{proof}

\end{document}
